\documentclass[12pt, a5paper]{article}

\usepackage{geometry}
\usepackage[utf8]{inputenc}
\usepackage[spanish]{babel}
\usepackage{amsmath, amssymb}
\usepackage{enumerate}
\usepackage{verbatim}
% \usepackage{tikz}
% \usetikzlibrary{babel}
\usepackage{multicol}
\usepackage[pdftex]{hyperref}
\hypersetup{pdftitle={Teoría Matricial},pdfauthor={Luis Carrillo Gutiérrez}}

\renewcommand{\familydefault}{\sfdefault}

\begin{document}
\section*{Matriz}
Es un \textbf{ordenamiento de coeficientes técnicos} (elementos) que pueden ser números, funciones, vectores, etc. dispuestos en \textit{filas} y \textit{columnas}, que sirven para plantear un sistema homogéneo de ecuaciones lineales que verifican ciertas reglas y propiedades del álgebra líneal. \\ 

\noindent Una matriz es una \textbf{disposición rectangular} de elementos distribuídos en filas y columnas. Los elementos en referencia, pueden ser números reales ($\mathbb{R}$), números complejos ($\mathbb{C}$), funciones o también matrices en condiciones especiales.

\subsection*{Orden de una matriz}
Una matriz $A$ tiene un orden según el número de elementos ordenados en $m$ filas y $n$ columnas. $A_{m \times n}$. Se dice que la matriz $A$ tiene orden $m$ por $n$. \\

\noindent Utilizaremos los corchetes para representar una \textbf{matriz}. Si el arreglo tiene $m$ filas y $n$ columnas diremos que la matriz es de orden $m \times n$. \\

\noindent {Ejemplo :} 
\noindent {\footnotesize La siguiente es una matriz de orden $3 \times 4$}
$$
\begin{bmatrix}
0 & -5 & 20 & 6 \\
1 & \;6 & 10 & 7 \\
3 &  \;7 & 1 & 2
\end{bmatrix}_{3 \times 4}
$$

\noindent {\footnotesize Una matriz de orden $3 \times 1$}
$$
\begin{bmatrix}
a \\
x^2 \\
y
\end{bmatrix}_{3 \times 1}
$$

\noindent {\footnotesize O una matriz de $2 \times 2$}
$$
\begin{bmatrix}
 8 & 4 \\
 -3 & -6
\end{bmatrix}_{2 \times 2}
$$

\noindent {\footnotesize La siguiente es una matriz (vector) de $1 \times 3$}
$$
\begin{bmatrix}
1 & 3 & -10
\end{bmatrix}_{1 \times 3}
$$

En forma concisa, una matriz de orden $m \times n$, se denota con un letra mayúscula, y sus elementos con letras minúsculas.

\subsection*{Elementos de una matriz}
\noindent Los elementos de una matriz se simboliza en forma general {\Large $a_{ij}$}, donde $i$ son las \textit{filas} y $j$ son las \textit{columnas}. Los mismos tienen variación $i = 1, 2, 3, 4, \ldots , m$ y $j = 1, 2, 3, 4, \ldots, n$. \\

\noindent Ejemplo:

% \begin{minipage}[t]{0.5\textwidth}
\begin{multicols}{2}
$A = \begin{bmatrix}
 2 & 5 \\
 4 & 8
\end{bmatrix}_{2 \times 2}$ \\ 
% \end{minipage}

%second column
% \begin{minipage}[t]{0.5\textwidth}
$B = \begin{bmatrix}
4 & 2 & 3 \\
5 & 3 & 2 \\
2 & 4 & 6
\end{bmatrix}_{3 \times 3}$
% \end{minipage}
\end{multicols}

\subsection*{Notación general de una matriz}
Tenemos elementos con ciertos sub-índices que nos indican la fila y columna donde se encuentra ubicado cada elemento, asimismo, nos indiquen cuantas filas y cuantas columnas tienen en total. \\

\noindent Sea la matriz $A$ de $m$ filas y $n$ columnas cuyos elementos son de la forma $a_{ij}$ indicándonos que es un elemento de la fila y de la columna definida.

$$
A = \begin{pmatrix}
a_{11} & a_{12} & a_{13} & \dots & a_{1n} \\
a_{21} & a_{22} & a_{23} & \dots & a_{2n} \\
a_{31} & a_{32} & a_{33} & \dots & a_{3n} \\
\vdots & \vdots & \vdots & \ddots & \vdots \\
a_{m1} & a_{m2} & a_{m3} & \dots & a_{mn} \\
\end{pmatrix}_{m \times n}
$$

\noindent También es usual decir : {\large $A_{m \times n} = [a_{ij}]_{m \times n}$}; donde $i = 1, 2, 3, \ldots, m$ y $j = 1, 2, 3, \ldots, n$ siendo {\Large $a_{ij}$} es el elemento que se encuentra en la \textbf{i-ésima} fila y la \textbf{j-ésima} columna.

\subsection*{Clasificación de matrices}
\noindent Las matrices se clasifican en : 
\begin{itemize}
\item{matriz rectangular}
\item{matriz cuadrada}
\end{itemize}

\subsubsection*{Matriz rectangular}
\noindent Se dice que una matriz es \textbf{rectangular} cuando el número de filas es diferentes al número de columnas, es decir $m \neq n$. \\

\noindent Ejemplo:

\begin{multicols}{2}
$C = \begin{bmatrix}
2 & 2 & 3 \\
4 & 1 & 2
\end{bmatrix}_{2 \times 3}$ \\

$D = \begin{bmatrix}
3 & 2 \\
2 & 5 \\
1 & 2
\end{bmatrix}_{3 \times 2}$
\end{multicols}

\subsubsection*{Vector}
\noindent Es una matriz que tiene solamente una columna (\textit{vector columna}) o una matriz que tiene sólo una fila (\textit{vector fila}).
\begin{itemize}
\item{Si $i = 1$ y $j = 1, 2, \ldots, n$ entonces {\large $A_{1 \times n} = [a_{ij}]_{1 \times n}$} es la llamada \textbf{matriz fila} o \textbf{vector fila} / Horizontal.
$$
E = \begin{vmatrix}2 & 1 & 3\end{vmatrix}_{1 \times 3}
$$
}
\item{Si $i = 1, 2, \ldots, n$ y $j = 1$ entonces {\large $A_{m \times 1} = [a_{ij}]_{m \times 1}$} es la llamada \textbf{matriz columna} o \textbf{vector columna} / Vertical
$$
F = \begin{pmatrix}
3 \\
2 \\
\frac1{2}\\
4 \end{pmatrix}_{4 \times 1}
$$}
\end{itemize}

\subsubsection*{Matriz cuadrada}
\noindent Se dice que una matriz es cuadrada, cuando el número de filas es igual al número de columnas; es decir $m = n$. \\

\noindent Ejemplo :

\begin{multicols}{2}
$G = \begin{vmatrix}
2 & 4 \\
3 & 5
\end{vmatrix}_{2 \times 2}$ \\

$H = \begin{vmatrix} 4 & 2 & 3 \\  5 & 3 & 2 \\ 2 & 4 & 6 \end{vmatrix}_{3 \times 3}$
\end{multicols}


\subsubsection*{Clasificación de matriz cuadrada}
La \textbf{matriz cuadrada} se puede clasificar de acuerdo a la utilidad que proporcione en la resolución de problemas o ecuaciones matemáticas. \\

% \noindent Para este caso específico partimos del siguiente ejemplo: Una empresa de confecciones producen camisas, sacos y pantalones a base de telas de los siguientes colores: blanco, negro y azul. Los agentes de ventas pueden representarse en una tabla de insumos como el siguiente:

% | Componentes |  | Productos|   |
% | --- | --- | --- | --- |
% | &nbsp; | Camisa | Saco | Pantalón |
% | tela blanca | 6 m | 6 m | 4 m |
% | tela negra | 12 m | 12 m | 10 m |
% | tela azul | 10 m | 8 m | 8 m |

\subsubsection*{Matriz del sistema ({\Large $A_s$})}
\noindent Es aquella matriz cuadrada formada por los coeficientes técnicos de insumo-producto del problema real. \\

\noindent Ejemplo:

$$
A_s = \begin{vmatrix}
6 & 6 & 4 \\
12 & 12 & 10 \\
10 & 8 & 8
\end{vmatrix}_{3 \times 3}
$$

\subsubsection*{Matriz transpuesta ({\Large ${A_s}^T$})}
Es aquella matriz cuadrada que se forma mediante la transposición (intercambio) de los elementos de filas en columnas o columnas en filas de la matriz del sistema.

\noindent Ejemplo:

$$
{A_s}^T = \begin{vmatrix} 6 & 12 & 10 \\
6 & 12 & 8 \\
4 & 10 & 8 \end{vmatrix}_{3 \times 3}
$$

\noindent Si {\large $A = [a_{ij}]_{m \times n}$}, es una matriz $m \times n$, entonces la \\ \texttt{TRANSPUESTA} de $A$, es la matriz {\large $A^t$ = $[b_{ij}]_{n \times m}$}, donde {\large $b_{ij} = a_{ji}$} para todo $i = 1, 2, 3, \ldots, n$ y $j = 1, 2, 3, \ldots, m$ \\

\noindent Ejemplo :

\begin{multicols}{2}
Si $A = \begin{vmatrix} 8 & 1 & 4 \\ -3 & 3 & 5 \\ 0 & \sqrt{2} & -1 \\ -6 & 4 & 2 \end{vmatrix}_{4 \times 3}$ \\

$A^t = \begin{vmatrix} 8 & -3 & 0 & -6 \\ 1 & 3 & \sqrt{2} & 4 \\ 4 & 5 & -1 & 2 \end{vmatrix}_{3 \times 4}$
\end{multicols}

\subsubsection*{Matriz adjunta ($Adj.\;{A_s}^T$)}
Es una matriz cuadrada que se forma mediante los cofactores que se obtienen de la matriz transpuesta y dichos cofactores tienen signos  alternados de positivo a negativo tanto en filas y columnas, finalmente se aplica determinantes para cada uno de los cofactores.


\subsubsection*{Cofactor}
Los cofactores se forman por los elementos libres que quedan al elegir un elemento específico de la matriz cuadrada y eliminando (omitiendo) al mismo tiempo su fila y columna correspondiente. El número de cofactores es igual al número de elementos de la matriz cuadrada referencial. \\

\noindent El cofactor se aplica a la matriz del sistema para calcular su determinante, como también a una matriz transpuesta para determinar su matriz adjunta.

\subsubsection*{Determinante}
\noindent El determinante es una función que aplicada a una matriz cuadrada que la transforma en un valor escalar o un número. \\
\\ El determinante se calcula como la sumatoria algebraica del producto entre los elementos de la primera fila por sus cofactores correspondientes.

\noindent Ejemplo :

Sea $A = \begin{vmatrix} 1 & 6\\ 2 & 3\end{vmatrix}_{2 \times 2}$ \\
la determinante de $A$ (se representa con $|A|$), se obtiene por: $|A| = 1 \cdot \begin{vmatrix} 3 \end{vmatrix} - 6 \cdot \begin{vmatrix} 2 \end{vmatrix}$

$$
|A_s| = 6 12 10 8 8 - 6
$$

\subsubsection*{Diagonales de una matriz cuadrada}
\noindent Una matriz cuadrada tiene dos clases de diagonales.
\begin{itemize}
\item{\textbf{Diagonal Principal} -- Son aquellas diagonales orientadas del vértice superior izquierdo hacia el vértice inferior derecho, es decir, de $a_{11}$ a $a_{mn}$ y tienen signo positivo.}
\item{\textbf{Diagonal Secundaria} -- Son aquellas diagonales orientadas del vértice superior derecho hacia el vértice inferior izquierdo , es decir, de $a_{1n}$ a $a_{m1}$ y tienen signo negativo.}
\end{itemize}

\subsubsection*{Matriz Inversa.}
Es el cociente que se obtiene al dividir la matrix adjunta y el determinante de la matriz del sistema. Es una matriz

\subsubsection*{Matriz unitaria o Identidad (I)}
\noindent Es aquella matriz cuya diagonal principal son unidades (1). \\

\noindent Ejemplo:
$$
I = \begin{vmatrix} 1 & 0 \\ 0 & 1\end{vmatrix}_{2 \times 2}
$$



\begin{comment}
```
np.empty((2,3))
```

## Igualdad de Matrices

Dos matrices $A$ y $B$ son iguales, si y sólo si, son del mismo

## Suma y resta de matrices
Para sumar o restar dos o más matrices, la condición necesaria es que sean del mismo orden, tanto matrices rectangulares y cuadradas respectivamente.

A = 1 2; 5 4

4 8

2 6

4 2 3
2 1 6
5 7 2

5 2 5
6 10 4
8 6 8

1 2
5 4

4 8
2 6

5 10
7 10


## Multiplicación de matrices
Para multiplicar dos o más matrices la condición necesaria será que el número de columnas del primer facto debe ser igual al número de filas del segundo factor.

El procedimiento para obtener

Ejemplo
Sean las siguientes matrices
As = [1 2 5 4 4 8 2 6

2 3 4
5 4 2
4 2 5
4 3 2]

Multiplicar
As Bs = 1 2 5 4

### Cálculo de la matriz inversa por el método matriz adjunta

Para calcular la matriz inversa, se debe identificar la matriz del sistema formados por los coeficientes técnicos de insumo-producto. Sólo las matrices cuadradas poseen la inversa.

El procedimiento es el siguiente:
PRIMERO
Se tiene la matriz del sistema siguiente:
Segundo
Se determina la matriz compuesta AT. Esta matriz se obtiene cambiando los elementos de la fila en columnas y los elemenos de la columna en filas.

tercero
Se calcula la matriz adjunta At está matriz adjunta se forma mediante los cofactores del mismo orden de la matriz transpuesta en la siguiente forma:

en una matriz transpuesta se elige el elemento de la primera fila y primera columna (a11), luego se anulan todos los elementos de la fila y columna que contengan dicho elemento.
\end{comment}
\end{document}